\section{Introduction}

Reclaiming space in an index requires reading all of it. If the pages are read in
an order known only as the traversal unfolds, each read must complete before the
next is issued, and the pass proceeds at one device round trip per page. If the
order is known in advance, several reads can be outstanding at once, and the pass
proceeds at the device's throughput rather than at its latency.

Maintenance of a generalized search tree acquired a known order when its
traversal was changed from logical to physical: pages are visited in increasing
block number. That change was made for locality~--- sequential rather than random
access~--- but it also makes the sequence predictable, which is the precondition
for issuing reads ahead of need. PostgreSQL~18 gained a streaming read interface
for exactly this, and we applied it to maintenance in three access methods.

This paper reports what happened when we measured it, which is nothing, and why
that is the correct outcome rather than an implementation failure. It also
reports the measurement that locates the missing gain: it was already being
collected by the operating system.

\paragraph{Contributions.}

\begin{enumerate}
\item A measurement showing no difference from application-level prefetching in
      index maintenance under default settings, with the explanation
      (Section~\ref{sec:nogain}).
\item A decomposition, by varying kernel readahead on one device, of how much of
      the benefit of ordered traversal is attributable to readahead rather than
      to locality as such (Section~\ref{sec:decompose}).
\item A statement of the conditions under which an application-level mechanism
      is nevertheless required (Section~\ref{sec:when}).
\end{enumerate}

\section{Background}

Maintenance visits every leaf and every internal page, removes entries pointing
at deleted rows, and reclaims pages that have become empty. Its cost is
%
\begin{equation}
\label{eq:vac}
P = (L + I)\,c_s + \beta\,c_r ,
\end{equation}
%
where $L$ and $I$ count leaf and internal pages, $c_s$ and $c_r$ are the costs of
a sequential and of a random page read, and $\beta$ counts pages that must be
read again because a concurrent insertion moved an unprocessed entry behind the
sweep. The first term is what ordered traversal buys; the second is what
concurrency costs, and its pages arrive in an order nobody knows in advance.

Prefetching acts on the first term in a way the notation hides. Issuing $d$ reads
concurrently does not reduce the number of reads; it replaces $n$ latencies by
$n/d$ of them, up to the point where the device's throughput binds. Whether this
is worth anything depends entirely on whether the reads were being issued one at
a time in the first place~--- and that is a property of the whole storage stack,
not of the database.

\section{The measurement that shows nothing}
\label{sec:nogain}

We compare two builds differing by exactly the commit that puts maintenance of a
generalized search tree onto the streaming read interface. Twenty million points,
every fifth row deleted, a buffer cache of 256~MB against an index of some
160\,000 pages, the operating system's page cache dropped before each
measurement, and the table warmed by a scan beforehand so that the only cold
pages are index pages.

Six runs give a median of 52.5~seconds on both sides. The two distributions are
not merely close; they are the same distribution, and the individual runs
interleave.

The reason is visible once one asks what the operating system was doing. The
default readahead on this device is 128~KB, that is sixteen pages: on detecting a
sequential access pattern the kernel fetches sixteen pages ahead of the reader.
Physical-order traversal produces exactly the pattern that triggers this. The
work that the application-level mechanism was written to do had therefore already
been done, one layer down, for free, and the second implementation of it adds
nothing.

We state this plainly because the temptation in reporting such work is to find
the configuration in which the number moves. That configuration exists, and
Section~\ref{sec:when} describes it; but the honest headline for a change shipped
to all users is what happens in the setting all users have.

\section{Locating the benefit by removing readahead}
\label{sec:decompose}

If the kernel is doing the work, one can measure how much work that is by
stopping it. We vary \code{read\_ahead\_kb} on one device and re-measure the
underlying comparison~--- logical against physical traversal order~--- on five
million points with the cache dropped and the table warmed.

\begin{table}[t]
\centering
\caption{Index cleanup time against readahead. Medians of three runs.}
\label{tab:ra}
\begin{tabular}{rrrr}
\toprule
\code{read\_ahead\_kb} & logical, s & physical, s & ratio \\
\midrule
0 & 57.8 & 52.8 & 1.09 \\
128 (default) & 20.9 & 6.1 & 3.43 \\
4096 & 16.7 & 6.2 & 2.69 \\
\bottomrule
\end{tabular}
\end{table}

With readahead disabled, ordered traversal is worth 9\,\%. With it enabled, 3.43
times. The entire advantage of knowing the page order is therefore realized
\emph{through} readahead: the ordering itself buys almost nothing directly,
because on this device a random read costs nearly what a sequential one does. The
number of buffer-cache misses differs by a factor of two in all three rows and is
unaffected by any of this, which confirms that what changes is not how many
requests are made but whether they can be merged and anticipated.

Two further observations from the same table. Physical-order traversal is
reproducible to 0.2~seconds while logical order is not, since the latter's cost
depends on how its scattered requests happen to fall. And a very large readahead
helps logical order rather than physical: at 4096~KB the former improves from
20.9 to 16.7 while the latter is unchanged at 6.2, already limited by bandwidth.
The measured ratio at 4096~KB is thus formally smaller than at the default, which
is a reminder that ``more readahead'' is not a monotone improvement.

\section{When the application-level mechanism is nevertheless needed}
\label{sec:when}

The negative result above is conditional on the kernel being able to help, and
there are several ways it cannot.

\paragraph{When there is no kernel in the path.} A database configured for direct
I/O bypasses the page cache and its readahead entirely. This is not exotic: it is
the configuration chosen when the database's own buffer cache is sized to the
machine and double buffering is judged wasteful.

\paragraph{When the pattern is not recognized.} Kernel readahead is a heuristic
over observed accesses. It recognizes a forward sequential scan of a file. The
$\beta$ term of~\eqref{eq:vac}~--- pages revisited because concurrent insertions
moved entries behind the sweep~--- is by construction not sequential, and no
heuristic can anticipate it. Those pages are read the ordinary way, and an
application-level mechanism that knew of them could not help either, since their
identity is discovered rather than predicted.

\paragraph{When the layer below has different information.} Network-attached and
virtualized storage may reorder, merge, or throttle requests according to rules
the kernel's heuristic does not model. Our own device is of this kind, which is
why its random and sequential costs are so close and why
Table~\ref{tab:ra} shows a smaller spread than a rotating disk would.

The correct claim for the change, then, is not that it makes maintenance faster.
It is that it makes maintenance's I/O pattern explicit rather than emergent, so
that performance no longer depends on a heuristic in another layer inferring what
the database already knows. That is worth having, and it is a different claim
from a speedup, and conflating the two is how performance work acquires its
reputation for unreproducibility.

\section{Limitations}

We measured one storage device, one workload and one access method in detail.
The negative result is therefore a statement about a common configuration rather
than about all configurations, and Section~\ref{sec:when} lists cases we have not
measured. In particular we have not measured direct I/O, which is the case where
we expect the mechanism to pay and where a positive result would be most
informative; we report its absence rather than presenting the argument as
complete.

The comparison of Section~\ref{sec:decompose} varies readahead on a single device
and therefore emulates only the removal of readahead, not the seek cost of a
rotating disk. An independent report on such a disk puts the advantage of ordered
traversal at fiftyfold~\cite{janes2019gistvacuum}, which we cannot reproduce and
do not treat as ours.

\section{Conclusion}

Application-level prefetching in index maintenance produced no measurable gain in
the configuration our users run, because the operating system was already
performing the same anticipation from a pattern the traversal order handed it for
free. Disabling readahead shows how much that was worth: the advantage of ordered
traversal falls from 3.43 to 1.09, meaning almost all of it was readahead and
almost none of it was locality as such on this class of device.

We report this as a result because the alternative~--- searching for the
configuration in which the number moves and presenting that~--- would describe a
system nobody runs. The mechanism's value is not speed in the default setting but
independence from a heuristic in another layer, which matters exactly where that
heuristic is absent or wrong. Establishing that positively requires measuring
those configurations, which we have not done, and we would rather leave the claim
incomplete than complete it with the wrong evidence.
